\documentclass{standalone}
\usepackage{circuitikz}
\usetikzlibrary{calc}
\definecolor{circuitnotemark}{HTML}{FE00FE}
\begin{document}
\begin{circuitikz}[american, line width=0.8pt]
\ctikzset{bipoles/length=1.2cm}
\ctikzset{grounds/scale=1.36}
\coordinate (b1) at (0,-2);
\coordinate (e1) at (0,-8);
\coordinate (b3) at (4,-2);
\coordinate (b5) at (8,-2);
\coordinate (b7) at (12,-2);
\coordinate (a3) at (4,0);
\coordinate (a5) at (8,0);
\coordinate (a7) at (12,0);
\coordinate (e7) at (12,-8);
\draw (b1) to[esource, l_=$\mathrm{W1}$] (e1); % line 3
\draw (-0.18,-5) sin ++(0.09,0.09) cos ++(0.09,-0.09) sin ++(0.09,-0.09) cos ++(0.09,0.09); % line 3
\node[ground] at (e1) {}; % line 4
\draw (b3) to[R, l_=$R_{1}$, a^=$1\,\mathrm{k}\Omega$] (b5); % line 5
\draw (b5) to[C, l_=$C_{1}$, a^=$100\,\mathrm{n}\mathrm{F}$] (b7); % line 6
\draw (a3) to[rmeter, t={$\mathrm{V}$}, l_=$\mathrm{CH1}$] (a5); % line 7
\draw (a5) to[rmeter, t={$\mathrm{V}$}, l_=$\mathrm{CH2}$] (a7); % line 8
\draw (b1) -- (b3); % line 10
\draw (a3) -- (b3); % line 11
\draw (a5) -- (b5); % line 12
\draw (a7) -- (b7); % line 13
\draw (b7) -- (e7); % line 14
\draw (e7) -- (e1); % line 15
\node[circ] at (e1) {};
\node[circ] at (b3) {};
\node[circ] at (b5) {};
\node[circ] at (b7) {};
\node[circ] at (a5) {};
\node[anchor=south west, inner sep=2pt, yshift=4pt, circuitnotemark, font=\LARGE\bfseries] (circuittitle) at (current bounding box.north west) {X};
\path (circuittitle.south west) rectangle ++(9.114,0.853);
\node[anchor=north east, inner sep=2pt, circuitnotemark, font=\large] (circuitstamp) at ([yshift=-0.593cm]current bounding box.south east) {X};
\path (circuitstamp.north east) rectangle ++(-5.08,-0.593);
\end{circuitikz}
\end{document}
